\scanpage{253}

% upper centre, folio
\mspage{6}

% upper left, square with diagonal
\diagram{\begin{tikzcd}[ampersand replacement=\&,column sep=small,row sep=small]
X\arrow[d,"p"']\&X\times_e Y\arrow[l,"q'"']\arrow[d,"p'"]\arrow[dl,dashed,"\varphi"']\\
e\&Y\arrow[l,"q"]
\end{tikzcd}}

% upper right, continuation from scan 252
s'établit en calculant le deuxième \unsure{mb|mbr} comme

% upper right, below continuation
\[
\varphi_!(M_{X\times Y})\simeq
p_!\bigl(\underbrace{q'_!(M_{X\times Y})}_{\substack{\rotatebox{90}{\(\simeq\)}\\q_!(M_Y)_X\\\text{car }q\text{ est}\\\mathbb D\text{-lisse}}}\bigr)
\simeq\unsure{\xi|\zeta}\times_{\mathbb D}(\unsure{\eta|y}\times_{\mathbb D}M)
\qquad\text{OK.}
\]

% before proposition
J'ai aussi envie de prouver

% proposition with interlinear hypothesis
\uline{Proposition !} \inserted{Soit $\xi\in\mathrm{Ob}\,\mathrm{Hot}_0$. Alors} $F\mapsto\xi\times_{\mathbb D}F$, en tant que morphisme de dérivateurs $\mathbb D\unsure{-|\to}\mathbb D$, est \uline{exact à droite}, et $F\mapsto\mathrm{hom}_{\mathbb D}(\xi,F)$ est exact à gauche.

% above middle display
\unsure{OBS} que $\xi$ est défini par $X$ dans $\mathrm{Diag}$. Soit $f:\deleted{\unsure{X|I}}I\longrightarrow J$, \deleted{\unsure{i|il}}soit $F$ dans $\mathbb D(I)$, il faut prouver

% middle, desired comparison
\[f_*(\mathrm{hom}_{\mathbb D}(\xi,F))\simeq \mathrm{Hom}_{\mathbb D}(\xi,f_*(F))\]

% above middle square
Considérons donc

% middle, comparison square
\diagram{\begin{tikzcd}[ampersand replacement=\&,column sep=large]
X\times I\arrow[r,no head,"f'"]\arrow[d,"p="',"{(x,i)\mapsto i}"]\& X\times J\arrow[d,"{q=(x,j)\mapsto j}"]\\
I\arrow[r,"f"]\&J
\end{tikzcd}}

% below square
On a

% first line below square
\[\mathrm{hom}_{\mathbb D}(\xi,F)=p_*(p^*F)\]

% between computations
donc

% composite direct images
\[f_*(\mathrm{hom}_{\mathbb D}(\xi,F))=\underbrace{f_*p_*}_{(fp)_*}(\unsure{y'|\gamma'|p'|p^*}F)\simeq(qf')_*p^*(F)\simeq q_*(f'_*p^*(F))\]

% above base-change formula
Mais

% base-change formula
\[f'_*p^*(F)\simeq q^*f_*(F)\]

% last three lines
formule de changemt de base des $f_*$, pour le chgt de base $q$, qui est lisse, car il dérive de $X\longrightarrow e$ qui est lisse par chgt. de base $J\longrightarrow e$. On trouve

% archive sheet number: 253
